\title{Direct Preference Optimization: Your Language Model is Secretly a Reward Model}

\begin{abstract}
Large-scale unsupervised language models (LMs) possess extensive world knowledge and exhibit some degree of reasoning capability, yet exerting fine-grained control over their outputs remains challenging because their training is entirely unsupervised. Common techniques to enhance their controllability rely on gathering human judgments about the comparative quality of model outputs and adapting the original unsupervised LM to reflect these human preferences, frequently leveraging reinforcement learning from human feedback (RLHF). Nonetheless, RLHF is known for its procedural complexity and instability, typically involving two steps: training a reward model to encode human preference data, and subsequently adjusting the large unsupervised LM using reinforcement learning in a way that optimizes the predicted reward while remaining close to the initial model. This paper presents a novel approach to parameterizing the reward model within RLHF, which permits closed-form derivation of the optimal corresponding policy. This reformulation allows the canonical RLHF challenge to be addressed using a straightforward classification loss. The method, termed \textit{Direct Preference Optimization} (DPO), is both efficient and robust, eliminating the requirements for LM sampling during fine-tuning and extensive hyperparameter search. Empirical results indicate that DPO can adapt LMs to better reflect human preferences at least as effectively as current leading approaches. In particular, DPO surpasses PPO-based RLHF methods in controlling sentiment in generated outputs, while offering comparable or superior performance in tasks such as summarization and single-turn dialogue—achieved with notable simplicity in both implementation and training procedures.
\end{abstract}

\section{Introduction}
Unsupervised language models (LMs) trained on massive corpora demonstrate a range of unexpected abilities~\citep{chowdhery2022palm, brown2020language, touvron2023llama, bubeck2023sparks}. Despite this, such models learn from human-generated data encompassing diverse intentions, levels of expertise, and objectives. Imitating all such behaviors is not always desirable; for instance, while an AI assistant for coding should \textit{recognize} typical programming errors to correct them, we ultimately want the assistant to \textit{emulate} the (possibly infrequent) expert-level solutions contained in the training data. Analogously, although it is useful for a language model to be \textit{informed} about widely held misconceptions (e.g., endorsed by 50\% of individuals), the goal is emphatically not for the model to propagate those misconceptions half the time they are queried. Thus, delineating and reliably inducing a model’s \emph{preferred outputs and conduct}—as distinct from its extensive \textit{capabilities and knowledge}—is essential to developing AI systems that are robust, effective, and aligned with user intentions~\citep{ouyang2022training}. While state-of-the-art approaches generally leverage reinforcement learning (RL) to tailor LMs to human preferences, we demonstrate that the RL-based objective commonly adopted can in fact be exactly optimized by a standard binary cross-entropy loss. This leads to a significant simplification of the preference learning workflow.

In broad terms, prevalent methods adapt LMs to exhibit targeted behaviors using carefully assembled datasets of human preferences that typify desired, safe, and useful outputs. This adaptation phase follows an initial, unsupervised large-scale pretraining over text. The most direct preference learning approach is to conduct supervised fine-tuning on human-produced high-quality responses; nevertheless, the predominant paradigm involves reinforcement learning from human or AI feedback (RLHF/RLAIF; \citep{christiano2017deep, bai2022constitutional}). RLHF first learns a reward model from human preference data, and then employs RL to train a model to favor responses with high reward while avoiding undue deviation from the pretrained model. While RLHF models exhibit remarkable performance in dialogue and code generation, the RLHF pipeline itself introduces additional layers of complexity relative to supervised learning. This stems from the necessity to maintain and train multiple LMs as well as to sample iteratively from the policy throughout training, which considerably raises computational demands.

In this work, we introduce a method for optimizing language models in alignment with human preferences, bypassing the need for explicit reward modeling or reinforcement learning frameworks. We present \textit{{Direct Preference Optimization} (DPO)}, a technique that implicitly targets the same objective as standard RLHF procedures (maximizing reward with a KL-divergence regularization), yet offers a more accessible and straightforward approach to implementation and training. Conceptually, the DPO update boosts the relative log likelihood of preferred answers over those deemed less favorable, employing a dynamic, per-sample importance weighting mechanism that helps to avert the model collapse observed with a simple probability ratio approach. In common with prior algorithms, DPO utilizes a theoretical model of preference, such as the Bradley-Terry framework \cite{bradley1952rankanalysis}, to assess the correspondence between a candidate reward function and observed human preference data. Distinctively, rather than leveraging this model to train an intermediary reward function which is then maximized by the policy, DPO applies a variable substitution to reformulate the preference loss directly in terms of the policy itself. Provided with annotated preference comparisons over generated responses, DPO enables the direct optimization of the policy via a binary cross entropy loss, yielding an optimal policy with respect to a latent reward derived from the empirical preferences.

The principal contribution of this study is the Direct Preference Optimization (DPO) method, which offers a straightforward, reinforcement learning–free framework for preference-driven language model training. Experimental results demonstrate that DPO matches or outperforms existing strategies, including PPO-based RLHF, in domains such as sentiment adaptation, summarization, and dialogue, utilizing language models scaling up to 6B parameters.

\section{Related Work}

As self-supervised language models scale up, they increasingly demonstrate the ability to handle certain tasks with zero-shot learning \citep{radford2019language} or through few-shot prompting strategies \citep{gpt3,megatron,chowdhery2022palm}. Nonetheless, their effectiveness on various downstream applications and their alignment with user intentions are markedly enhanced when fine-tuned using datasets comprising explicit instructions and corresponding human-generated completions \citep{mishra-etal-2022-cross,sanh2022multitask,chung2022scaling,thoppilan2022lamda}. The process of 'instruction-tuning' not only allows LLMs to adapt to instructions that were absent from the fine-tuning data but also generally bolsters their practical usability \citep{chung2022scaling}. Although instruction-tuning has achieved impressive results, collecting \textit{relative} human assessments of response quality is often more straightforward than obtaining high-quality expert demonstrations. Consequently, later studies have utilized datasets reflecting human preferences to further fine-tune LLMs, which has led to gains in areas such as translation \citep{kreutzer-etal-2018-reliability}, summarization \citep{stiennon2022learning,ziegler2020finetuning}, narrative generation \citep{ziegler2020finetuning}, and adherence to instructions \citep{ouyang2022training,ramamurthy2023is}. These methodologies typically begin by training a neural reward model to align with a preference dataset, commonly using the Bradley-Terry framework \citep{bradley1952rankanalysis}. Subsequently, a language model is further fine-tuned to optimize for this reward via reinforcement learning techniques such as REINFORCE \citep{williams1992reinforce}, proximal policy optimization (PPO; \cite{schulman2017proximal}), or related approaches \citep{ramamurthy2023is}. Closely related research explores the generation of synthetic preference datasets emphasizing specific features—like safety or harmlessness—using LLMs refined for instruction-following via human feedback, with only weak human oversight (e.g., rubrics to guide LLM annotation) \citep{bai2022constitutional}. This approach intersects two significant research threads: the use of reinforcement learning to train language models on diverse objectives~\citep{Ranzato2015SequenceLT,paulus2018a,wu2018learning}, and the broader methodologies for learning from human preference data \citep{christiano2017deep,kupcsik2018learning}. However, in spite of the theoretical advantages of leveraging comparative human preference data, applying reinforcement learning at scale to large language models is still logistically difficult. This study introduces a principled alternative for optimizing language models with respect to relative preferences that circumvents the need for reinforcement learning.

Beyond language tasks, learning policies based on preferences has been investigated in both the bandit and reinforcement learning communities, resulting in the development of multiple methods. In contextual bandit settings, when preferences or rankings over actions (rather than explicit reward signals) are provided, this is referred to as the contextual dueling bandit (CDB; \cite{yue2012karmed,dudik2015contextual}). Without access to absolute reward values, the theoretical framework of CDBs replaces the goal of identifying an optimal policy with the concept of a \textit{von Neumann winner}: a policy that secures an expected victory against \textit{every} other candidate policy at least 50\% of the time \citep{dudik2015contextual}. Notably, CDB algorithms usually operate with preference feedback delivered in an online manner. In contrast, research on learning from human preferences frequently works with a predetermined, offline collection of action pairs that have already been annotated with preferences \citep{yan2022human}. Likewise, \textit{preference-based RL} (PbRL) methods are driven by binary preference feedback generated via an \textit{unknown} `scoring' function as opposed to direct reward feedback \citep{BusaFekete2014,ruiz2023dueling}. While a variety of algorithms have been proposed for PbRL, including those leveraging off-policy preference datasets, they typically entail explicitly learning a latent reward or scoring function before optimizing the policy according to it \citep{jain2013learning,BusaFekete2014,christiano2017deep,sadigh2017active,kupcsik2018learning}. In contrast, we introduce a unified approach that trains a policy in a single stage to align directly with preference signals.

\section{Preliminaries}\label{section:prelims}

We summarize the RLHF methodology as described in \citeauthor{ziegler2020finetuning} and subsequently adopted in \citep{stiennon2022learning, bai2022training, ouyang2022training}. This pipeline is usually divided into three main components: (1) supervised fine-tuning (SFT); (2) preference elicitation and reward modelling; and (3) reinforcement learning optimization.

\textbf{SFT}: The RLHF workflow starts by adapting a large pre-trained language model to the desired application (such as conversation or summarization) using supervised learning on a curated, high-quality dataset, thereby obtaining the supervised fine-tuned model $\pi^\text{SFT}$.

\textbf{Reward Modelling Phase}: During this stage, prompts $x$ are fed into the supervised model to generate answer pairs $(y_1, y_2)\sim \pi^\text{SFT}(y \mid x)$. Human evaluators are then tasked with expressing a preference between the two responses, denoted as $y_w\succ y_l \mid x$ where $y_w$ and $y_l$ correspond to the preferred and non-preferred responses, respectively. These choices are presumed to be influenced by some unobserved reward function $r^*(y, x)$ to which we do not have direct access. A range of preference modelling approaches exist, with the Bradley-Terry (BT) model \cite{bradley1952rankanalysis} commonly employed; the more general Plackett-Luce model \citep{plackett1975analysis, luce2012individual} is also applicable when several ranked options are available. In the BT setting, the underlying distribution over human preferences $p^*$ is expressed as follows:

\begin{equation}\label{eq:bradley-terry}
    p^*(y_1\succ y_2 \mid x)=\frac{\exp\left(r^*(x, y_1)\right)}{\exp\left(r^*(x, y_1)\right) + \exp\left(r^*(x, y_2)\right)}.
\end{equation}

Given a fixed dataset of comparison samples $\mathcal{D}=\bigl\{x^{(i)}, y_w^{(i)}, y_l^{(i)}\bigr\}_{i=1}^N$ drawn from the target distribution $p^*$, we can introduce a parameterized reward function $r_{\phi}(x, y)$ and determine the optimal parameters by maximizing the likelihood of the data. By recasting this task as binary classification, the loss to minimize becomes the negative log-likelihood:

\begin{equation}\label{eq:reward_model}
    \mathcal{L}_R(r_{\phi}, \mathcal{D}) = -\mathbb{E}_{(x, y_w, y_l)\sim \mathcal{D}}\bigl[\log \sigma(r_{\phi}(x, y_w)- r_{\phi}(x, y_l))\bigr]
\end{equation}

with $\sigma$ denoting the logistic sigmoid. For language models, $r_{\phi}(x, y)$ is commonly constructed by augmenting the final transformer block of the SFT model $\pi^\text{SFT}(y \mid x)$ with a linear projection, yielding a single real-valued reward per input-output pair~\cite{ziegler2020finetuning}. To promote stability in the learned reward, standard practices include normalizing outputs to ensure $\mathbb{E}_{x,y\sim \mathcal{D}}\left[r_\phi(x, y)\right] = 0$ across $x$.

\textbf{RL Fine-Tuning Phase}: In the reinforcement learning step, this trained reward function is employed as the signal for further tuning the language model. In line with earlier literature~\citep{jaques2017sequence, jaques2020human}, the following optimization is adopted:

\begin{equation}\label{eq:RL}
\max_{\pi_{\theta}}  \mathbb{E}_{x\sim \mathcal{D}, y\sim \pi_{\theta}(y \mid x)}\bigl[r_{\phi}(x, y)\bigr] - \beta\mathbb{D}_{\textrm{KL}}\bigl[\pi_{\theta}(y\mid x)\mid \mid \pi_\text{ref}(y\mid x)\bigr],
\end{equation}

where $\beta$ modulates the regularization towards a base reference policy $\pi_\text{ref}$, specifically the SFT initialization $\pi^\text{SFT}$. The policy $\pi_\theta$ is similarly initialized to the SFT parameters in actual implementations. This constraint on divergence is crucial for restricting the language model's outputs to regions where the reward estimator is reliable, maintaining diverse outputs, and avoiding collapse to a limited set of maximally scored responses. Since sampling discrete tokens from the model makes this objective non-differentiable, optimization is performed using reinforcement learning methods. The typical procedure~\cite{ziegler2020finetuning, stiennon2022learning, bai2022training, ouyang2022training} is to define the combined reward ${r(x, y) = r_{\phi}(x, y) -\beta (\log \pi_{\theta}(y\mid x) - \log \pi_\text{ref}(y\mid x))}$ and employ PPO~\cite{schulman2017proximal} to optimize the objective.

\section{Direct Preference Optimization}\label{sec:DPO}

Recognizing the complexity and computational demands of applying reinforcement learning techniques for fine-tuning large language models, our objective is to design a streamlined policy optimization method that operates directly on preference data. Departing from conventional RLHF pipelines, which first learn a reward surrogate before optimizing it with RL, we exploit a special form of reward model parameterization that admits a closed-form solution for its optimal policy—thereby eliminating the need for iterative RL optimization. The core idea, which we elaborate below, is to use an explicit analytical relationship between a reward function and its corresponding optimal policy, facilitating a transformation from loss functions on rewards to loss functions over policies. This variable-substitution approach circumvents the need to fit an explicit reward model, yet preserves fidelity to established frameworks for modeling human preferences such as the Bradley-Terry model. As a result, the policy network encapsulates both the generative

\textbf{Derivation of the DPO Objective.} Our starting point aligns with prior research, utilizing the RL objective detailed in Eq.~\ref{eq:RL} for a general reward function $r$. Building on previous findings~\citep{peters2007reinforcement, peng2019advantage, korbak2022reinforcement, go2023aligning}, it follows directly that the policy maximizing reward while constrained by KL-divergence, as posed in Eq.~\ref{eq:RL}, adopts the structure:

\begin{equation}\label{eq:op_policy}
    \pi_r(y\mid x) = \frac{1}{Z(x)}\pi_\text{ref}(y\mid x)\exp\left(\frac{1}{\beta}r(x, y)\right),
\end{equation}

where $Z(x) =\sum_{y}\pi_\text{ref}(y\mid x)\exp\left(\frac{1}{\beta}r(x, y)\right)$ denotes the partition function. A comprehensive derivation is available in Appendix \ref{app:derivation1}. However, even with a maximum likelihood estimate $r_{\phi}$ of the true reward $r^*$, evaluating $Z(x)$ remains computationally intensive~\citep{korbak2022reinforcement, go2023aligning}, presenting practical challenges for directly employing this form. Nevertheless, one may re-express Eq.~\ref{eq:op_policy} to solve for the reward in terms of its optimal policy $\pi_r$, the reference $\pi_\text{ref}$, and the yet-unknown $Z(\cdot)$. Specifically, applying a logarithm to both sides and rearranging terms leads to:

\begin{equation}\label{eq:main_eq}
    r(x,y) =\beta \log \frac{\pi_r(y\mid x)}{\pi_\text{ref}(y\mid x)} + \beta \log Z(x).
\end{equation}

Applying this formulation to the true reward $r^*$ and its optimal policy $\pi^*$, we observe that the Bradley-Terry model is determined solely by reward differences between pairs of completions: ${p^*(y_1 \succ y_2 \mid x) = \sigma(r^*(x, y_1) - r^*(x, y_2))}$. Substituting the above parameterization of $r^*(x,y)$ into the preference model from Eq.~\ref{eq:bradley-terry}, the partition function terms are eliminated, allowing human preference probabilities to be written entirely with respect to the optimal policy $\pi^*$ and the reference policy $\pi_\text{ref}$. Thus, the optimal RLHF policy $\pi^*$ for the Bradley-Terry model fulfills the preference relationship:

\begin{equation}\label{eq:objective}
    p^*(y_1\succ y_2 \mid x)=\frac{1}{1 + \exp\left(\beta \log \frac{\pi^*(y_2\mid x)}{\pi_\text{ref}(y_2\mid x)} - \beta \log \frac{\pi^*(y_1\mid x)}{\pi_\text{ref}(y_1\mid x)}\right)}
\end{equation}

Details of this derivation are provided in Appendix~\ref{app:derivation2}. While Eq.~\ref{eq:objective} is based on the Bradley-Terry assumption, analogous developments are feasible for the broader Plackett-Luce family~\citep{plackett1975analysis, luce2012individual}, as detailed in Appendix~\ref{app:plackett_luce_models}.

Having related the likelihood of observed human preferences directly to the optimal policy (bypassing the explicit reward representation), we can cast a maximum likelihood estimation problem for a policy model $\pi_\theta$. In analogy to the reward modeling setup (see Eq.~\ref{eq:reward_model}), the resulting policy optimization objective is:

\begin{equation}\label{eq:optimum_model}
    \mathcal{L}_\text{DPO}(\pi_{\theta}; \pi_\text{ref}) = -\mathbb{E}_{(x, y_w, y_l)\sim \mathcal{D}}\left[\log \sigma \left(\beta \log \frac{\pi_{\theta}(y_w\mid x)}{\pi_\text{ref}(y_w\mid x)} - \beta \log \frac{\pi_{\theta}(y_l\mid x)}{\pi_\text{ref}(y_l\mid x)}\right)\right].
\end{

In this framework, we infer an implicit reward by employing an alternative parameterization in which the optimal policy coincides with $\pi_\theta$. Additionally, because our method aligns with optimizing a reparameterized Bradley-Terry model, it inherits favorable theoretical characteristics, including consistency under appropriate assumptions on the distribution of preference data \cite{bong2022generalized}. A more detailed discussion of these theoretical aspects and how they relate to related literature is presented in Section~\ref{sec:theory}.

\textbf{Mechanistic insights into the DPO update.} To gain an operational perspective on DPO, we examine the gradient of the loss $\mathcal{L}_\text{DPO}$. The derivative with respect to $\theta$ is given by:

\begin{multline*}\label{eq:gradient}
    \nabla_\theta \mathcal{L}_\text{DPO}(\pi_\theta;\pi_\text{ref}) = \\ -\beta\mathbb{E}_{(x, y_w, y_l) \sim \mathcal{D}} \bigg[\underbrace{\sigma(\hat{r}_\theta(x, y_l) - \hat{r}_\theta (x, y_w))}_\text{higher weight when reward estimate is wrong}\bigg[\underbrace{\nabla_\theta\log \pi(y_w \mid x)}_\text{increase likelihood of $y_w$} - \underbrace{\nabla_\theta\log\pi(y_l \mid x)}_\text{decrease likelihood of $y_l$}\bigg]\bigg],
\end{multline*}

Here, $\hat{r}_\theta(x, y) = \beta \log \frac{\pi_\theta(y \mid x)}{\pi_\text{ref}(y \mid x)}$ denotes the implicit reward derived from the language model $\pi_\theta$ relative to the reference $\pi_\text{ref}$ (see Section~\ref{sec:theory} for further elaboration). Conceptually, the loss gradient for $\mathcal{L}_\text{DPO}$ drives up the likelihood of preferred continuations $y_w$ and lowers the likelihood for less preferred alternatives $y_l$. Crucially, each sample’s influence is modulated by how much the implicit reward model $\hat{r}_\theta$ overrates the dispreferred output compared to the preferred one, amplified by the parameter $\beta$—thereby reflecting the extent to which the implicit reward function misranks outputs and factoring in the KL regularization strength. Empirical observations highlight that this weighting is critical: omitting the weighting factor can induce pathological model behavior, as evidenced in Appendix Table~\ref{tab:unlikelihood_generations}.

\textbf{DPO overview.} The DPO process typically involves the following steps: 1) For each prompt $x$, generate completion samples $y_1, y_2 \sim \pi_\text{ref}(\cdot \mid x)$ and obtain human preference labels to assemble the offline preference dataset $\mathcal{D} = \{x^{(i)}, y_w^{(i)}, y_l^{(i)}\}_{i=1}^N$, and 2) train the language model $\pi_\theta$ to minimize $\mathcal{L}_\text{DPO}$, conditioned on $\pi_\text{ref}$, $\mathcal{D}$, and the selected $\beta$. 
In real-world scenarios, it is often preferable to leverage publicly available preference datasets instead of creating completions and acquiring new human annotations. When such preference data is obtained via $\pi^\text{SFT}$, we set $\pi_\text{ref} = \pi^\text{SFT}$ if it exists. In instances where $\pi^\text{SFT}$ is unavailable, $\pi_\text{ref}$ is initialized by maximizing the likelihood over the preferred completions ${(x, y_w)}$, specifically, by solving ${\pi_\text{ref} = \argmax_{\pi}\mathbb{E}_{x, y_w \sim \mathcal{D}}\left[\log \pi(y_w \mid x)\right]}$. This initialization aims to address the gap between the inaccessible true reference distribution and the $\pi_\text{ref}$ utilized by DPO. Additional information regarding implementation details and hyperparameter selection can be found in Appendix~\ref{app:implementation}.

\section{Theoretical Analysis of DPO} This section provides a deeper examination of the DPO technique, delivering theoretical foundations and highlighting the benefits of DPO in comparison with actor-critic approaches for RLHF, such as PPO~\cite{schulman2017proximal}.

\label{sec:theory}

\subsection{Your Language Model Is Secretly a Reward Model} DPO achieves policy learning in a single maximum likelihood step, thereby eliminating the need for explicit reward modeling or separate RL-based policy optimization. Specifically, the optimization problem in Eq. \ref{eq:main_eq} corresponds to the Bradley-Terry model with rewards specified as $r^*(x, y) = \beta \log\frac{\pi^*_\theta(y \mid x)}{\pi_\text{ref}(y \mid x)}$. The parameterized model $\pi_{\theta}$ is optimized analogously to the reward model in Eq. \ref{eq:reward_model}, given an appropriate change of variables. In this discussion, we develop the theoretical framework supporting this reparameterization, demonstrate that it does not limit the reward models that can be learned, and prove that it enables exact recovery of the optimal policy. We begin by formalizing an equivalence relation for reward functions.

\begin{definition}
Two reward functions $r(x, y)$ and $r'(x, y)$ are defined to be equivalent if and only if there exists a function $f$ such that ${r(x, y) - r'(x, y) = f(x)}$.
\end{definition}

It can be readily verified that this definition satisfies the conditions of an equivalence relation, leading to a partitioning of all reward functions into equivalence classes. This gives rise to the following two lemmas:

\begin{lemma}\label{lemma:same_prefrence}
Within the framework of Plackett-Luce models (including the Bradley-Terry model as a special case), reward functions that belong to the same equivalence class generate identical preference distributions.
\end{lemma}

\begin{lemma}\label{lemma:same_policy}
    Any two reward functions from the same equivalence class result in the same optimal policy under a constrained RL setting.
\end{lemma}

The demonstrations of these lemmas are elementary and can be found in Appendix \ref{app:lemma1}. The first lemma formalizes a standard under-identification property of the Plackett-Luce family \cite{plackett1975analysis}. This non-identifiability necessitates the introduction of additional constraints to ensure identifiability and obtain reliable MLE estimates as in Eq. \ref{eq:reward_model} \cite{bong2022generalized}. According to the second lemma, every reward function from an equivalence class leads to the same optimal policy, implying that for our objectives, it is sufficient to recover any representative reward function from the optimal equivalence class. The following Theorem is proven in Appendix~\ref{app:thm1}:

\begin{theorem}\label{thm:main}
    Under suitable mild conditions, all equivalence classes of reward functions compatible with the Plackett-Luce (and particularly Bradley-Terry) models are expressible through the reparameterization ${r(x, y) = \beta \log \frac{\pi(y\mid x)}{\pi_\text{ref}(y\mid x)}}$ for some model $\pi(y\mid x)$ and chosen reference policy $\pi_\text{ref}(y \mid x)$.
\end{theorem}

\begin{sproof}
    Take any reward function $r(x, y)$ that gives rise to its associated optimal policy $\pi_r(y \mid x)$ as specified in Eq. \ref{eq:op_policy}. We will demonstrate that a representative of the equivalence class containing $r$ can be written via the above reparameterization. Define the projection operator $f$ by
\begin{equation}
    f(r; \pi_\text{ref}, \beta)(x, y) = r(x, y) - \beta\log\sum_{y}\pi_\text{ref}(y\mid x)\exp\left(\frac{1}{\beta}r(x, y)\right)
\end{equation}
Here, $f$ adjusts the reward function by subtracting the log-partition function with respect to the reference policy $\pi_\text{ref}$, evaluated over all possible $y$. Since this normalization only depends on $x$, the adjusted reward $f(r; \pi_\text{ref}, \beta)(x, y)$ remains in the equivalence class of $r(x, y)$. If we now substitute $r$ with the right-hand side of Eq.~\ref{eq:main_eq} (valid for any reward), we obtain $f(r; \pi_\text{ref}, \beta)(x, y) = \beta \log \frac{\pi_r(y\mid x)}{\pi_\text{ref}(y\mid x)}$. This establishes that the projection $f$ produces a representative in the required form, and the proposed reparameterization does not reduce the generality of the reward representation.
\end{sproof}

Theorem~\ref{thm:main} can also be interpreted as identifying the specific reward function chosen by the DPO reparameterization within each equivalence class: namely, the reward function that fulfills the condition

\begin{equation}\label{eq:lag_p}
     \sum_{y}\underbrace{\pi_\text{ref}(y\mid x)\exp\left(\frac{1}{\beta}r(x, y)\right)}_{=\pi(y\mid x)\text{, using Thm.~\ref{thm:main} reparam.}} = 1,
\end{equation}

meaning that $\pi(y\mid x)$ defines a legitimate probability distribution, with all probabilities being non-negative and summing to unity. Examining Eq.~\ref{eq:lag_p} in light of Eq.~\ref{eq:op_policy}, it becomes evident that this expression corresponds to the partition function of the optimal policy associated with the reward function $r(x, y)$. The central idea of the DPO method lies in strategically constraining the generally under-determined Plackett-Luce (specifically, Bradley-Terry) class of preference models. By doing so, the range of representable reward models remains intact, while ensuring that the optimal policy described by Eq.~\ref{eq:op_policy} can be computed in closed-form for every prompt $x$.

\subsection{Instability of Actor-Critic Algorithms} Our framework is equally useful for analyzing sources of instability within standard actor-critic approaches in RLHF settings, such as PPO. Restricting attention to the RL fine-tuning phase (Section \ref{section:prelims}), we relate this procedure to the control-as-inference paradigm \cite{levine2018reinforcement} in the context of the constrained RL formulation of \ref{eq:RL}. Assuming a parameterized policy $\pi_{\theta}(y\mid x)$, the objective is to minimize the KL divergence $\mathbb{D}_{\text{KL}}[\pi_{\theta}(y|x) \mid \mid \pi^*(y\mid x)]$, where $\pi^*$ denotes the optimal policy obtained in Eq. \ref{eq:optimum_model} for a given reward $r_{\phi}(y, x)$. This setup yields, after appropriate manipulation, the following optimization criterion:

\begin{equation}\label{eq:AC}
    \max_{\pi_{\theta}}\mathbb{E}_{\pi_{\theta}(y\mid x)}\bigg[\underbrace{r_{\phi}(x, y) -\beta\log\sum_{y}\pi_\text{ref}(y\mid x)\exp\left(\frac{1}{\beta}r_{\phi}(x, y)\right)}_{f(r_{\phi}, \pi_\text{ref}, \beta)} - \underbrace{\beta\log\frac{\pi_{\theta}(y\mid x)}{\pi_\text{ref}(y\mid x)}}_{\text{KL}}\bigg]
\end{equation}

This objective has been optimized in earlier studies 
\citep{ziegler2020finetuning, stiennon2022learning, bai2022training, ouyang2022training} using the DPO-equivalent reward formulation within the reward class $r_{\phi}$. In this framework, the normalization component in $f(r_{\phi}, \pi_\text{ref}, \beta)$ can be interpreted as the soft value function associated with the reference policy $\pi_\text{ref}$. Although this normalization does not alter the set of optimal policies, omitting it can result in a policy gradient with large variance, thereby introducing instability into the training process. Incorporating the normalization term through a parameterized value function is possible, yet this approach is often challenging to train effectively. Another strategy found in the literature uses a human-generated completion as a baseline, effectively providing a Monte Carlo estimate of the normalization term from a single instance. By contrast, the DPO reparameterization generates a reward structure that eliminates the need for any such baseline.

\section{Experiments}
This section presents an empirical analysis of DPO’s capacity for policy learning directly from human preferences. We first investigate, within a controlled text-generation environment, how effectively DPO balances reward maximization against KL-divergence regularization relative to widely used preference-based algorithms like PPO. Subsequently, we examine DPO’s applicability to more substantial models and challenging RLHF benchmarks, including tasks in summarization and conversational agents. Our results indicate that, even with minimal hyperparameter adjustment, DPO consistently matches or outperforms competitive methods such as RLHF with PPO and surpasses selecting the top trajectory among $N$ samples as scored by a learned reward model. Prior to detailing these findings, we outline the experimental design; for further information, see Appendix~\ref{app:exp_details}.

\textbf{Tasks.} We assess algorithmic performance on three distinct open-ended text generation challenges. Across all experiments, the objective is to derive a policy from a dataset of preferences $\mathcal{D}=\bigl\{x^{(i)}, y_w^{(i)}, y_l^{(i)}\bigr\}_{i=1}^N$. For \textbf{controlled sentiment generation}, each $x$ serves as the initial segment of an IMDb movie review \cite{maas-EtAl:2011:ACL-HLT2011}, and the policy is tasked with generating a continuation $y$ exhibiting positive sentiment. To maintain experimental control, we create preference pairs for model outputs by applying a pre-trained sentiment classifier, ensuring that $p(\text{positive}\mid x,y_w) > p(\text{positive}\mid x,y_l)$. The SFT approach involves fine-tuning GPT-2-large to convergence using movie reviews from the IMDB training split (see App~\ref{app:sentiment_details} for additional methodology). For the \textbf{summarization} task, $x$ represents a Reddit forum post, and the policy must compose a concise summary $y$ encapsulating the core ideas. Consistent with earlier studies, we utilize the Reddit TL;DR summarization dataset \citep{volske-etal-2017-tl}, supplemented with human preference annotations from \citeauthor{stiennon2022learning}. Here, we rely on an SFT model adapted to human-authored summaries of forum posts\footnote{\url{https://huggingface.co/CarperAI/openai_summarize_tldr_sft}}, implemented using the TRLX \citep{leandro_von_werra_2023_7790115} toolkit for RLHF. The preference data collected by \citeauthor{stiennon2022learning} are derived from outputs of a comparable, albeit independently fine-tuned, SFT model. The third task, \textbf{single-turn dialogue}, frames $x$ as a user query spanning a broad range of topics from science to personal matters. Here, the policy is responsible for generating a reply $y$ that is both constructive and engaging. For this purpose, we utilize the Anthropic Helpful and Harmless dialogue dataset \citep{bai2022training}, comprising 170,000 human-assistant exchanges, where each record concludes with a pair of model-generated responses accompanied by a label indicating the human-preferred answer. In this scenario, we lack access to a ready-made SFT model; instead, we perform fine-tuning on a standard language model using only preferred responses to establish the SFT baseline.

\textbf{Evaluation.} We employ two distinct evaluation strategies in our experiments. To rigorously compare how well each algorithm optimizes the objective of constrained reward maximization, we focus on the controlled sentiment generation task by measuring the trade-off between the attained reward and KL-divergence from the reference policy. The full Pareto frontier can be precisely computed in this setting, since we have access to the ground-truth reward metric (provided by a sentiment classifier). However, in real-world scenarios where the true reward function is unknown, we assess each algorithm’s \textit{win rate} compared to a baseline, using GPT-4 as a stand-in for human evaluation. In the summarization task, the baseline comprises reference summaries from the test split, whereas for single-turn dialogue, we use the human-preferred responses as the baseline. While previous work indicates that language models can outperform existing metrics as automated evaluators \citep{Chen2023ExploringTU}, we corroborate our reliance on GPT-4 through a supplementary human evaluation study (see Sec.~\ref{sec:human-judgments}). Our findings demonstrate that GPT-4’s ratings align closely with human judgment, with the rate of human agreement with GPT-4 comparable to or even surpassing the agreement among different human annotators.

\textbf{Methods.} Alongside DPO, we assess multiple established techniques for steering language models towards alignment with human preferences. At the simplest level, we consider zero-shot prompting of \textbf{GPT-J} \citep{gpt-j} for summarization and 2-shot prompting with \textbf{Pythia-2.8B} \citep{biderman2023pythia} for dialogue. Further, we include the \textbf{SFT} model as well as \textbf{Preferred-FT}, a model trained via supervised learning using only the preferred completions $y_w$—these are generated either by the SFT model (in controlled sentiment and summarization) or a base LM (for dialogue tasks). We also experiment with the pseudo-supervised \textbf{Unlikelihood}~\citep{welleck2019neural} approach, which trains models to assign higher probability to $y_w$ and lower probability to $y_l$; an adjustable coefficient $\alpha\in[0,1]$ controls the strength of the unlikelihood penalty. Our study also encompasses \textbf{PPO} \citep{schulman2017proximal} leveraging a reward function learned from human preference data and \textbf{PPO-GT}, an oracle variant that utilizes the true reward function available in controlled sentiment settings. For sentiment tasks, two implementations of PPO-GT are tested: one is an out-of-the-box solution \cite{leandro_von_werra_2023_7790115} and the other incorporates reward normalization and additional hyperparameter tuning for enhanced results (these enhancements are also applied to PPO runs using learned rewards). As a final strong baseline, we evaluate the \textbf{Best of $N$} method, which samples $N$ candidate responses from the SFT (or Preferred-FT for dialogue), then selects the one rated highest by a learned reward model. While this decouples reward modeling from PPO optimization and tends to yield strong results, it is computationally demanding since it requires $N$ completions per input during inference.

\subsection{How well can DPO optimize the RLHF objective?}

The standard KL-regularized reward maximization formulation central to RLHF methodologies ensures that policies not only maximize the reward but also remain close to the reference policy, constraining over-exploration. Thus, evaluations of algorithmic performance must jointly consider both the achieved reward and the corresponding KL divergence; securing marginally greater rewards at the cost of substantial increases in KL is suboptimal. Figure~\ref{fig:frontier-tldr-main} presents the trade-off curve between reward and KL divergence for several approaches within the sentiment analysis context. For each algorithm, we perform multiple independent runs, varying a policy conservativeness parameter each time (target KL $\in\{3,6,9,12\}$ for PPO, $\beta \in \{0.05,0.1,1,5\}$, $\alpha\in\{0.05,0.1,0.5,1\}$ for unlikelihood, and different random seeds for preferred-FT), yielding a total of 22 runs. Every 100 training steps up to convergence, each trained policy is assessed over a set of evaluation prompts, measuring both mean reward according to the actual reward function and mean sequence-level KL\footnote{This denotes the aggregate sum of KL-divergences calculated at each timestep.} relative to the reference policy, denoted as $\text{KL}\left(\pi\mid \mid \pi_\text{ref}\right)$. The findings demonstrate that DPO consistently yields a more favorable frontier, achieving superior rewards with comparatively modest KL divergence. This is especially significant for several reasons. Firstly, both DPO and PPO aim to solve the same underlying optimization problem, yet DPO attains this outcome with greater efficiency—its reward/KL relationship universally surpasses PPO. Secondly, DPO's frontier remains superior to PPO even when PPO is provided access to exact ground truth rewards (PPO-GT).

\subsection{Can DPO scale to real preference datasets?}
\label{sec:dpo-real-datasets}
We proceed by assessing how well DPO fine-tunes on tasks involving summarization and single-turn dialogue. In the context of summarization, standard automatic metrics such as ROUGE often exhibit weak alignment with human judgment~\citep{stiennon2022learning}. Earlier research has shown that optimizing LMs using PPO with human preference data results in higher-quality summaries. To benchmark the various approaches, we sample model outputs from the test split of the TL;DR summarization dataset and calculate the mean win rate when compared against ground truth references. Sampling is conducted across a range of temperatures between 0.0 and 1.0, with win rates reported in Figure~\ref{fig:frontier-tldr-main} (right). The DPO, PPO, and Preferred-FT models all start from the same GPT-J SFT checkpoint\footnote{\url{https://huggingface.co/CarperAI/openai_summarize_tldr_sft}}. Our experiments reveal that DPO achieves a win rate close to 61\% at a temperature of 0.0, outperforming PPO, which records about 57\% at its optimal temperature of 0.0. Notably, DPO attains a higher peak win rate than the $N$ baseline method. It is important to point out that DPO’s $\beta$ hyperparameter was not subject to extensive tuning in these trials, suggesting these findings may not reflect its full capacity. Additionally, DPO demonstrates considerably greater stability across temperature variations in comparison to PPO, whose performance can drop to baseline GPT-J levels at higher temperatures. The Preferred-FT approach, in contrast, offers little improvement beyond the SFT model. Finally, a direct comparison between DPO and PPO in human evaluations, detailed in Section~\ref{sec:human-judgments}, shows that DPO samples at a temperature of 0.25 were chosen 58\% of the time over PPO outputs generated at the same temperature.

For single-turn dialogue evaluation, we analyze various approaches using the portion of the Anthropic HH dataset’s test split \citep{bai2022training} consisting of single exchanges between a human and an assistant. In assessments involving GPT-4, we employ the dataset’s preferred completions as references, enabling calculation of win rates for each approach. As there is no standardized SFT baseline for this scenario, our process begins with a pre-trained Pythia-2.8B model; we apply Preferred-FT to finetune a reference model on selected completions to ensure the model’s outputs remain in-distribution, and subsequently, we train the model with DPO. For benchmarking, we include the strongest result among 128 Preferred-FT completions (noting that the Best of $N$ method does not surpass performance at $N=128$, as shown in Appendix Figure~\ref{fig:best-of-n}) and test a 2-shot prompted configuration of the original Pythia-2.8B model. We observe that DPO matches or exceeds the best-performing temperature settings for these baselines. Additionally, we assess an RLHF model that uses PPO and has been trained on the Anthropic HH dataset \footnote{\url{https://huggingface.co/reciprocate/ppo_hh_pythia-6B}}, referencing an established repository \footnote{\url{https://github.com/CarperAI/trlx/tree/main/examples/hh}}, but we are unable to identify a prompt or sampling temperature that yields results superior to the unmodified Pythia-2.8B model. Considering our findings on TL;DR and the shared reward optimization objectives of both methods, we use Best of 128 as an approximate indicator of PPO-level performance. In summary, DPO stands out as the sole computationally efficient strategy that consistently surpasses the preferred completions in the Anthropic HH dataset, achieving equal or improved outcomes relative to the more resource-intensive Best of 128 method. Moreover, Figure~\ref{fig:dialogue-main} illustrates that DPO achieves near-optimal results after relatively few iterations.

\subsection{Generalization to a new input distribution}

To further investigate how PPO and DPO perform when faced with shifts in data distribution, we assess the summarization policies learned from our Reddit TL;DR experiment on an out-of-distribution benchmark: the test split of the CNN/DailyMail dataset, which consists of news articles~\citep{nallapati-etal-2016-abstractive}. We retain the optimal sampling temperatures identified for the TL;DR setting (0 and 0.25, respectively). Table~\ref{tab:ood} summarizes these results. We used GPT-4 to measure win rates versus gold-standard summaries by adapting the original Reddit TL;DR comparison prompt—substituting “forum post” for “news article.” On this novel data domain, DPO continues to surpass the PPO-derived policy by a notable margin. These findings suggest that DPO can generalize to unfamiliar data as effectively as PPO, even though PPO benefits from access to supplementary, unlabeled Reddit TL;DR prompts that DPO does not exploit.

\subsection{Validating GPT-4 judgments with human judgments}
\label{sec:human-judgments}
We carry out a human evaluation to test the validity of GPT-4's preference judgments, drawing on outputs from the TL;DR summarization study and employing two distinct prompts for GPT-4. The \textbf{GPT-4 (S)} (simple) prompt directly queries which summary better conveys the salient information from the source post. In contrast, the \textbf{GPT-4 (C)} (concise) prompt incorporates an additional criterion of conciseness; we consider this prompt because with the simple variant, GPT-4 often rates more verbose, repetitive outputs higher than humans do. Appendix~\ref{app:prompts} contains the complete prompts used. For diversity in summary quality, we conduct three sets of comparisons: the best DPO model (temperature 0.25), the weakest PPO model (temperature 1.0), and an intermediate SFT model (temperature 0.25). Each is pitted against PPO at its optimal temperature via greedy decoding. Analysis reveals that, for both prompt styles, GPT-4's preferences correlate with human judgments to a degree similar to inter-human agreement, indicating GPT-4 can serve as a reasonable substitute for direct human evaluation (note that, given our limited pool of human annotators, multiple ratings per item were obtained only for DPO and PPO-1 comparisons). In general, win rates from the \textbf{GPT-4 (C)} prompt most closely match human choices; hence, this prompt is selected for reporting main results in Section~\ref{sec:dpo-real-datasets}. Details about the experimental protocol, such as the rater interface and the identities of human participants, are provided in Appendix~\ref{app:human-study}.

\section{Discussion}
Learning from preferences offers an effective and scalable methodology for developing robust, aligned language models. In this work, we presented DPO, a straightforward approach that enables the training of language models from preference data without the reliance on reinforcement learning. Unlike approaches that reshape the preference learning objective to fit conventional RL frameworks and leverage standard RL techniques, DPO establishes an explicit correspondence between reward functions and language model policies. This permits training to align with human preferences \textit{directly}, via a simple cross-entropy loss, foregoing the complexity of reinforcement learning and without compromising expressivity. DPO achieves results on par with, or superior to, leading RLHF methods—such as those leveraging PPO—with minimal hyperparameter tuning, thereby substantially lowering the barrier to aligning language models using human feedback.

\textbf{Limitations \& Future Work.} Our findings prompt several avenues for further inquiry. What is the out-of-distribution generalization capacity of DPO policies, especially relative to methods that utilize explicit reward functions? Preliminary evidence indicates that DPO matches PPO-based approaches in this regard, but more thorough analysis is required. Another open question is whether training on self-generated labels from the DPO policy could harness unlabeled prompts effectively. Additionally, the manifestation of reward over-optimization in direct preference optimization warrants investigation, including whether the modest decline observed in Figure~\ref{fig:dialogue-main}-right constitutes such a phenomenon. While our evaluation covers models up to 6B parameters, scaling DPO to much larger, state-of-the-art architectures represents an exciting opportunity for future exploration. With respect to evaluation, our results show that GPT-4-based win rates can depend heavily on prompt phrasing, indicating that further research into optimal strategies for eliciting reliable automated judgments is necessary. Lastly, the utility of DPO extends beyond aligning language models with human preferences and may encompass applications such as training generative models in different domains.